\documentclass[11pt]{article}

\usepackage[T1]{fontenc}
\usepackage[margin=1in]{geometry}
\usepackage{amsmath,amssymb,amsthm}
\usepackage{graphicx}
\usepackage{hyperref}
\usepackage{enumitem}

\hypersetup{
  colorlinks=true,
  linkcolor=blue,
  urlcolor=blue
}

\newcommand{\R}{\mathbb{R}}
\newcommand{\E}{\mathbb{E}}
\newcommand{\Prob}{\mathbb{P}}
\newcommand{\fl}{\operatorname{fl}}
\newcommand{\Frob}[1]{\left\lVert #1\right\rVert_F}
\newcommand{\Norm}[1]{\left\lVert #1\right\rVert_2}
\newcommand{\Atil}{\widetilde A}
\newcommand{\Util}{\widetilde U}
\newcommand{\Ahat}{\widehat{\widetilde A}}
\newcommand{\Uhat}{\widehat{\widetilde U}}
\newcommand{\Ghat}{\widehat G}
\newcommand{\D}{D}
\newcommand{\uunit}{u}
\newcommand{\gammas}{\gamma_s}
\newcommand{\OneMat}{\mathbf 1_{n\times n}}
\newcommand{\lean}[1]{\texttt{\detokenize{#1}}}
\newcommand{\leanfit}[1]{\noindent{\scriptsize\nolinkurl{#1}}\par}

\newtheorem{theorem}{Theorem}
\newtheorem{corollary}[theorem]{Corollary}
\newtheorem{lemma}[theorem]{Lemma}
\theoremstyle{definition}
\newtheorem{definition}[theorem]{Definition}
\theoremstyle{remark}
\newtheorem*{remark}{Remark}
\newenvironment{formalized}
  {\begin{quote}\small\noindent\textbf{Lean formalization.}\par\smallskip}
  {\end{quote}}

\title{Formal Stability Bounds for RandNLA Algorithm 2}
\author{LeanFpAnalysis repository}
\date{}

\begin{document}
\sloppy

\maketitle

\begin{abstract}
This note states, in standard mathematical form, the Lean formalization of
Algorithm 2 from Drineas and Mahoney's CACM RandNLA survey,
\href{https://dl.acm.org/doi/10.1145/2842602}{RandNLA: Randomized Numerical
Linear Algebra}.  The first part records the row-sampling Frobenius guarantee
for the squared-row-norm distribution.  The second part specializes the same
formal machinery to leverage-score sampling, equation (6), and proves the
subspace-embedding estimate of equation (7), including the additional
floating-point perturbation term.
\end{abstract}

\section{Notation and Algorithm}

Let \(A\in\R^{m\times n}\).  For its \(i\)-th row write \(A_{i*}\), and set
\[
  r_i(A)=\Norm{A_{i*}}^2=\sum_{j=1}^n A_{ij}^2,\qquad
  \D(A)=\Frob{A}^2=\sum_{i=1}^m r_i(A).
\]
Algorithm 2 samples indices \(I_1,\ldots,I_s\) independently from a row
distribution \(p=(p_1,\ldots,p_m)\).  The sampled and rescaled sketch is
\[
  \Atil_{tj}=\frac{A_{I_tj}}{\sqrt{s\,p_{I_t}}},
  \qquad 1\le t\le s,\quad 1\le j\le n .
\]
The exact sampled Gram matrix is
\[
  G(\Atil)=\Atil^T\Atil,\qquad
  G(A)=A^TA.
\]

The floating-point model has unit roundoff \(\uunit\).  The repository uses
\(\gamma_s\) for the standard dot-product error factor, with the Lean
side-condition \(\lean{gammaValid fp s}\).  We distinguish two computed Gram
matrices:
\[
  \widehat G_{\rm scale}=G(\Ahat),
  \qquad
  \widehat G_{\rm dot}=\fl_{\rm dot}(\Ahat)^T\fl_{\rm dot}(\Ahat),
\]
where \(\Ahat\) denotes the sketch obtained after rounded row scaling, and
\(\widehat G_{\rm dot}\) computes each Gram entry using the repository's
floating-point dot-product algorithm.

\section{Equation (4): Squared-Row-Norm Sampling}

\begin{definition}[Squared-row-norm probabilities]
Assume \(\D(A)>0\).  The equation (4) distribution is
\[
  p_i=\frac{r_i(A)}{\D(A)}
  =
  \frac{\sum_j A_{ij}^2}{\sum_{k,\ell} A_{k\ell}^2}.
\]
\end{definition}

\begin{lemma}[Probability and support]
If \(\D(A)\ne0\), then \(\sum_i p_i=1\).  If \(\D(A)>0\), then
\(p_i\ge0\) for all \(i\).  Moreover, if row \(i\) contains a nonzero entry,
then \(p_i>0\), and for \(s>0\) the scaling denominator
\(\sqrt{s\,p_i}\) is nonzero.
\end{lemma}

\begin{formalized}
\lean{rowSqNormProb_sum_eq_one}\par
\lean{rowSqNormProb_nonneg}\par
\lean{rowSqNormProb_pos_of_entry_ne_zero}\par
\lean{rowSampleScaleDen_ne_zero}
\end{formalized}

\begin{lemma}[One rounded sampled row entry]
For a sampled row with \(p_{I_t}>0\), the rounded scaling operation satisfies
\[
  \left|\widehat{\Atil}_{tj}-\Atil_{tj}\right|
  \le
  \uunit\,|\Atil_{tj}|.
\]
\end{lemma}

\begin{formalized}
\lean{fl_rowSampleSketch_error_bound}.
\end{formalized}

\begin{remark}
Algorithm 2 does not aggregate repeated samples into a single row.  Repeated
draws produce repeated output rows.  Consequently the relevant random object is
the sampled Gram matrix \(\Atil^T\Atil\), not a grouped hit-count update.
\end{remark}

\section{Exact Gram Approximation}

\begin{theorem}[Unbiased sampled Gram entries]
Assume \(\D(A)>0\) and \(s>0\).  Under the independent row sampler with
probabilities \(p_i=r_i(A)/\D(A)\),
\[
  \E\bigl[(\Atil^T\Atil)_{jk}\bigr]=(A^TA)_{jk}
  \qquad\text{for all }j,k.
\]
\end{theorem}

\begin{formalized}
\leanfit{rowSqNormTraceProbability_expectationReal_rowSampleGram_entry}
\end{formalized}

\begin{theorem}[Second moment for the Frobenius error]
Assume \(\D(A)>0\) and \(s>0\).  Then
\[
  \E\!\left[\Frob{\Atil^T\Atil-A^TA}^2\right]
  \le
  \frac{\D(A)^2}{s}.
\]
\end{theorem}

\begin{formalized}
\leanfit{rowSqNormTraceProbability_expectationReal_sampleAverage_sub_mean_sq}

\leanfit{rowSqNormTraceProbability_expectationReal_rowSampleGram_frob_error_sq_le}
\end{formalized}

\begin{theorem}[Equation (5)]
Assume \(\D(A)>0\), \(s>0\), and \(\varepsilon>0\).  Then
\[
  \Prob\!\left[
    \Frob{\Atil^T\Atil-A^TA}
    \le
    \varepsilon\,\D(A)
  \right]
  \ge
  1-\frac{1}{s\varepsilon^2}.
\]
Equivalently, since \(\D(A)=\Frob{A}^2\),
\[
  \Prob\!\left[
    \Frob{\Atil^T\Atil-A^TA}
    \le
    \varepsilon\,\Frob{A}^2
  \right]
  \ge
  1-\frac{1}{s\varepsilon^2}.
\]
\end{theorem}

\begin{formalized}
\leanfit{FiniteProbability.eventProb_le_ge_one_sub_expectationReal_sq_div}

\leanfit{rowSqNormTraceProbability_eventProb_rowSampleGram_frob_error_le_epsilon}
\end{formalized}

\section{Floating-Point Perturbation of Equation (5)}

The formalized perturbation bounds are deterministic on the probability-one
support where every sampled row has positive sampling probability.  Define
\(\OneMat\) to be the \(n\times n\) all-ones matrix.  The Lean budgets are
Frobenius norms of constant \(n\times n\) matrices:
\[
  \tau_{\rm scale}(A)
  =
  \Frob{(2\uunit+\uunit^2)\D(A)\OneMat},
\]
\[
  \tau_{\rm dot}(A)
  =
  \Frob{\gammas(1+\uunit)^2\D(A)\OneMat},
  \qquad
  \tau_{\rm full}(A)=\tau_{\rm scale}(A)+\tau_{\rm dot}(A).
\]
Equivalently, since these scalar entries are nonnegative,
\[
  \tau_{\rm scale}(A)=n(2\uunit+\uunit^2)\D(A),
  \qquad
  \tau_{\rm dot}(A)=n\gammas(1+\uunit)^2\D(A).
\]
These are worst-case Frobenius budgets for, respectively, rounded row scaling
and rounded dot products in the Gram computation.

\begin{theorem}[Rounded scaling, exact Gram summation]
Assume \(\D(A)>0\), \(s>0\), and \(\varepsilon>0\).  If the sketch entries are
rounded but \(G(\Ahat)=\Ahat^T\Ahat\) is formed as an exact real Gram matrix,
then
\[
  \Prob\!\left[
    \Frob{\widehat G_{\rm scale}-A^TA}
    \le
    \varepsilon\,\D(A)+\tau_{\rm scale}(A)
  \right]
  \ge
  1-\frac{1}{s\varepsilon^2}.
\]
\end{theorem}

\begin{formalized}
\lean{rowSampleGramFpPerturbBudget}\par
\lean{rowSampleGramFpPerturbBudget_eq_nat_mul}\par
\lean{rowSampleGram_perturb_budget_le_explicit}

\leanfit{rowSqNormTraceProbability_eventProb_fl_rowSampleGram_frob_error_le_epsilon_add_explicit_budget}
\end{formalized}

\begin{theorem}[Fully floating-point Gram approximation]
Assume \(\D(A)>0\), \(s>0\), \(\varepsilon>0\), and
\(\lean{gammaValid fp s}\).  If each Gram entry is computed by the repository's
floating-point dot-product algorithm after Algorithm 2 has already returned its
sampled-and-scaled rows, then
\[
  \Prob\!\left[
    \Frob{\widehat G_{\rm dot}-A^TA}
    \le
    \varepsilon\,\D(A)+\tau_{\rm full}(A)
  \right]
  \ge
  1-\frac{1}{s\varepsilon^2}.
\]
\end{theorem}

\begin{formalized}
\lean{dotProduct_error_bound}\par
\lean{rowSampleGramDotProductBudget}\par
\lean{rowSampleGramDotProductBudget_eq_nat_mul}\par
\lean{rowSampleGramFullFpPerturbBudget}

\leanfit{rowSqNormTraceProbability_eventProb_fl_rowSampleGramDot_frob_error_le_epsilon_add_explicit_budget}
\end{formalized}

\begin{corollary}[\(\tau_{\rm dot}=0\) scaling-only specialization]
Assume \(\D(A)>0\), \(s>0\), and \(\varepsilon>0\).  If Algorithm 2 only
computes the sampled-and-scaled rows in floating point, and the Gram matrix is
introduced afterwards as the exact mathematical matrix
\(\widehat G_{\rm scale}=\Ahat^T\Ahat\), then the appropriate analysis model has
zero dot-product budget:
\[
  \tau_{\rm dot}(A)=0.
\]
Consequently,
\[
  \Prob\!\left[
    \Frob{\widehat G_{\rm scale}-A^TA}
    \le
    \varepsilon\,\D(A)+\tau_{\rm scale}(A)
  \right]
  \ge
  1-\frac{1}{s\varepsilon^2},
\]
where
\[
  \tau_{\rm scale}(A)=n(2\uunit+\uunit^2)\D(A).
\]
This is formalized as a separate theorem about the theoretical Gram matrix
formed from rounded rows, not as the computed-dot-product matrix with an
ignored dot-product error.  It is the faithful Algorithm 2 floating-point
stability statement when the Gram matrix appears only in the analysis.
\end{corollary}

\begin{formalized}
\lean{rowSampleGramFpPerturbBudget_eq_nat_mul}\par
\leanfit{rowSqNormTraceProbability_eventProb_fl_rowSampleGram_frob_error_le_epsilon_add_scaling_budget}
\leanfit{rowSqNormTraceProbability_eventProb_fl_rowSampleGram_frob_error_le_epsilon_add_tau_dot_zero}
\end{formalized}

\section{Equation (6): Leverage-Score Sampling}

Let \(U\in\R^{m\times n}\) have orthonormal columns:
\[
  U^TU=I_n.
\]
The \(i\)-th leverage score is
\[
  \ell_i=\Norm{U_{i*}}^2.
\]
Since
\[
  \sum_{i=1}^m \ell_i
  =
  \Frob{U}^2
  =
  \operatorname{trace}(U^TU)
  =
  n,
\]
equation (6) is
\[
  p_i
  =
  \frac{\ell_i}{\sum_r \ell_r}
  =
  \frac{\ell_i}{n}.
\]

\begin{theorem}[Formal equation (6)]
If \(U^TU=I_n\), then the squared-row-norm distribution applied to \(U\)
satisfies
\[
  p_i=\frac{\Norm{U_{i*}}^2}{n}
  \qquad\text{and}\qquad
  \sum_i p_i=1
\]
for \(n>0\).
\end{theorem}

\begin{formalized}
\lean{HasOrthonormalColumns}\par
\lean{leverageScore}\par
\lean{leverageScoreProb}\par
\leanfit{rowSqNormProbDen_eq_nat_of_orthonormal_columns}

\leanfit{leverageScoreProb_eq_rowNormSq_div_nat}

\lean{leverageScoreProb_sum_eq_one}
\end{formalized}

\section{Equation (7): Subspace Embedding}

The repository states the operator-2 estimate in vector-action form.  For a
square matrix \(M\), define
\[
  \operatorname{Op2Le}(M,c)
  \quad\Longleftrightarrow\quad
  \forall x\in\R^n,\ \Norm{Mx}\le c\Norm{x}.
\]
The formal proof uses the standard implication
\[
  \Frob{M}\le c
  \quad\Longrightarrow\quad
  \operatorname{Op2Le}(M,c).
\]

\begin{formalized}
\lean{vecNorm2}\par
\lean{opNorm2Le}\par
\leanfit{vecNorm2_matMulVec_le_frobNorm_mul}

\lean{opNorm2Le_of_frobNorm_le}
\end{formalized}

\begin{theorem}[Equation (7), exact arithmetic]
Let \(U\in\R^{m\times n}\) satisfy \(U^TU=I_n\).  Assume \(n>0\), \(s>0\),
and \(\varepsilon>0\).  Let \(\Util\) be the Algorithm 2 sketch produced by
sampling rows of \(U\) with the leverage probabilities \(p_i=\Norm{U_{i*}}^2/n\).
Then
\[
  \Prob\!\left[
    \forall x\in\R^n,
    \Norm{(\Util^T\Util-I_n)x}
    \le
    \varepsilon\Norm{x}
  \right]
  \ge
  1-\frac{1}{s(\varepsilon/n)^2}.
\]
\end{theorem}

\begin{formalized}
\lean{rowGram_eq_id_of_orthonormal_columns}.

\leanfit{leverageTraceProbability_eventProb_rowSampleGram_opNorm2_error_le_epsilon}
\end{formalized}

\begin{corollary}[Equation (7), concrete computed-denominator floating-point arithmetic]
Assume the exact Bennett sample-budget hypotheses for equation (7), and assume
\(\lean{gammaValid fp s}\).  Let \(\widehat G_{\rm dot}(U)\) be the
floating-point sampled Gram matrix formed from the exact leverage law
\(p_i=\|U_{i*}\|_2^2/n\), the concrete denominator routine
\[
  \widehat d_i=\operatorname{flsqrt}(\operatorname{flmul}(s,p_i)),
\]
rounded row divisions, and rounded length-\(s\) Gram dot products.  Then the
two-sided finite-Loewner event
\[
  \widehat G_{\rm dot}(U)-I_n
  \preceq
  \bigl(\varepsilon+\tau_{\rm comp}(U;s)\bigr)I_n,
  \qquad
  I_n-\widehat G_{\rm dot}(U)
  \preceq
  \bigl(\varepsilon+\tau_{\rm comp}(U;s)\bigr)I_n
\]
has probability at least \(1-\delta\), with
\[
  \tau_{\rm comp}(U;s)
  =
  n^2\left((2\eta_U+\eta_U^2)+\gamma_s(1+\eta_U)^2\right).
\]
Here \(\eta_U\) is the effective row-scaling relative error induced by the
same computed denominator routine.  The matrix \(U\) is an exact supplied
orthonormal-column basis; this corollary does not compute \(U\) from a data
matrix.
\end{corollary}

\begin{formalized}
\lean{leverage_fl_rowSampleGramDot_perturb_bound}.

\leanfit{leverageTraceProbability_eventProb_fl_rowSampleGramDot_opNorm2_error_le_epsilon_add_budget}
\leanfit{leverageFlMulThenSqrtRowScaleDen}
\leanfit{leverageTraceProbability_eventProb_fl_rowSampleGramDotWithFlMulThenSqrtDen_two_sided_finiteLoewnerLe_ge_one_sub_delta_of_sample_budget}
\end{formalized}

\begin{corollary}[Equation (7), actual input \(A=UC\) and computed Gram]
Let \(U\in\R^{m\times r}\) satisfy \(U^TU=I_r\), let
\(C\in\R^{r\times q}\), and set \(A=UC\).  The exact leverage law is
\(p_i=\|U_{i*}\|_2^2/r\), and the implementation samples rows of the actual
matrix \(A\).  The exact analysis witnesses \(U,C\) are not computed by this
theorem.

Under the same Bennett sample-budget hypotheses with \(r\) in place of \(n\),
the exact sampled actual-input Gram satisfies
\[
  \Pr\!\left[
    \widetilde A_U^T\widetilde A_U-A^TA\preceq \varepsilon A^TA
    \ \hbox{and}\
    A^TA-\widetilde A_U^T\widetilde A_U\preceq \varepsilon A^TA
  \right]\ge 1-\delta .
\]
With the concrete denominator
\(\widehat d_i=\fl_{\rm sqrt}(\fl_{\rm mul}(s,p_i))\), rounded row divisions
on entries of \(A\), and rounded length-\(s\) Gram dot products, the computed
Gram satisfies
\[
  \Pr\!\left[
    \widehat G_{\rm dot}(A;U)-A^TA
      \preceq \varepsilon A^TA+T_A^{\rm fp}I_q
    \ \hbox{and}\
    A^TA-\widehat G_{\rm dot}(A;U)
      \preceq \varepsilon A^TA+T_A^{\rm fp}I_q
  \right]\ge 1-\delta .
\]
Here
\[
  T_A^{\rm fp}
  =
  \left(\gamma_s(1+\eta_U^A)^2+2\eta_U^A+(\eta_U^A)^2\right)
  \sum_{j=1}^q
  \left(\sum_{i:\,p_i>0}\frac{|A_{ij}|}{\sqrt{p_i}}\right)^2,
\]
where
\[
  \eta_U^A
  =
  u+(1+u)
  \sum_{i:\,p_i>0}
  \frac{\sqrt{s p_i}\,(\sqrt{1+u}\,u+u)}
       {|\fl_{\rm sqrt}(\fl_{\rm mul}(s,p_i))|}.
\]
For fixed \(U,C,m,r,q,s\), this radius is
\[
  \Theta\!\left(
    (\gamma_s+u(1+|\{i:p_i>0\}|))
    \sum_{j=1}^q
    \left(\sum_{i:\,p_i>0}\frac{|A_{ij}|}{\sqrt{p_i}}\right)^2
  \right)
\]
and tends to zero as \(u,\gamma_s\to0\) when positive-probability computed
denominators remain nonzero.
\end{corollary}

\begin{formalized}
\leanfit{leverageFactoredInputSampleGram}
\leanfit{leverageTraceProbability_eventProb_factoredInputSampleGram_two_sided_finiteLoewnerLe_ge_one_sub_delta_of_sample_budget}
\leanfit{leverageFactoredInputDenGramBudget}
\leanfit{leverage_fl_rowSampleGramDotWithComputedDen_factoredInput_perturb_bound}
\leanfit{leverageTraceProbability_eventProb_factoredInput_fl_rowSampleGramDotWithFlMulThenSqrtDen_two_sided_finiteLoewnerLe_ge_one_sub_delta_of_sample_budget}
\end{formalized}

\begin{corollary}[Equation (7), stored basis table and computed Gram]
Let \(U\in\R^{m\times n}\) satisfy \(U^TU=I_n\), and let the exact leverage
law be \(p_i=\|U_{i*}\|_2^2/n\).  The probabilities and samples are exact
mathematical objects.  The three closed stored-basis modes are
\[
  \widehat U^{(\mathrm{mul1})}_{ij}=\fl_{\rm mul}(U_{ij},1),\qquad
  \widehat U^{(\mathrm{add0})}_{ij}=\fl_{\rm add}(U_{ij},0),\qquad
  \widehat U^{(\mathrm{sub0})}_{ij}=\fl_{\rm sub}(U_{ij},0),
\]
and all satisfy
\[
  |\widehat U^{(\mu)}_{ij}-U_{ij}|\le u|U_{ij}|.
\]
With
\[
  \rho_U=
  \sum_{i:\,p_i>0}
  \frac{\sqrt{s p_i}\,(\sqrt{1+u}\,u+u)}
       {|\fl_{\rm sqrt}(\fl_{\rm mul}(s,p_i))|},
  \qquad
  \lambda_U=2u+u^2+(1+u)^2\rho_U ,
\]
define the expanded deterministic radius
\[
  T_U^{\rm fp}
  =
  \left(\gamma_s(1+\lambda_U)^2+2\lambda_U+\lambda_U^2\right)
  \sum_{j=1}^n
  \left(\sum_{i:\,p_i>0}\frac{|U_{ij}|}{\sqrt{p_i}}\right)^2 .
\]
Under the exact Bennett sample-budget hypotheses and \(\gamma_s\) validity,
the Gram matrix computed from \(\widehat U^{(\mu)}\), the concrete denominator
\(\widehat d_i=\fl_{\rm sqrt}(\fl_{\rm mul}(s,p_i))\), rounded row divisions,
and rounded length-\(s\) dot products satisfies
\[
  \Pr\!\left[
    \widehat G_{\rm dot}^{(\mu)}-I_n\preceq
      (\varepsilon+T_U^{\rm fp})I_n
    \ \hbox{and}\
    I_n-\widehat G_{\rm dot}^{(\mu)}\preceq
      (\varepsilon+T_U^{\rm fp})I_n
  \right]\ge 1-\delta
\]
for every \(\mu\in\{\mathrm{mul1},\mathrm{add0},\mathrm{sub0}\}\).
This charges the stored basis table, denominator formation, row divisions, and
Gram dot products.  It does not claim that a QR or SVD routine has computed
\(U\) from a data matrix.
\end{corollary}

\begin{formalized}
\leanfit{RowSamplingLeverageComputedBasis.lean}
\leanfit{leverageComputedBasisDenGramBudget}
\noindent Stored-basis final wrappers are provided for the mul-one, add-zero,
and sub-zero modes.
\end{formalized}

\section{Hypotheses and Scope}

\begin{itemize}[leftmargin=2em]
  \item The exact equation (7) theorem assumes only \(U^TU=I_n\), \(n>0\),
  \(s>0\), and \(\varepsilon>0\).
  \item The source-aligned finite-Loewner equation (7) theorem assumes the
  displayed Bennett sample-budget inequalities and \(n>1\).
  \item The floating-point equation (7) corollaries additionally assume
  \(\lean{gammaValid fp s}\), inherited from the dot-product theorem.
  \item The older vector-action/Frobenius-to-operator consequence remains
  available, but the final equation (7) theorem surface is the formalized
  Oliveira/Tropp-style finite-Loewner Bennett route.
  \item The actual-input theorem uses exact analysis witnesses \(A=UC\) and
  computes with rows of \(A\); it charges the concrete denominator, row
  divisions, and Gram dot products.
  \item The stored-basis theorem charges the modeled rounded-copy table used
  in the sketch.  It still does not formalize a QR, SVD, or other routine that
  generates \(U\) from a data matrix.
  \item The proof does not formalize computation or approximation of leverage
  probabilities.  Equation (6) is the exact input sampling distribution once an
  orthonormal-column matrix \(U\) is fixed.
  \item The proof uses the vector-action operator predicate
  \(\operatorname{Op2Le}\) rather than introducing a supremum-valued spectral
  norm object.
\end{itemize}

\section{File Map}

\begin{itemize}[leftmargin=2em]
  \item \lean{RowSampling.lean}: row probabilities, product traces, sampled
  row scaling, and local floating-point scaling stability.
  \item \lean{RowSamplingGram.lean}: Gram matrices, unbiasedness, second
  moments, equation (5), and the floating-point Gram perturbation theorems.
  \item \lean{RowSamplingLeverage.lean}: equation (6) and equation (7) for
  leverage-score row sampling.
  \item \lean{RowSamplingLeverageMGF.lean}: source-aligned Bennett
  finite-Loewner equation (7), including the concrete computed-denominator
  floating-point endpoint.
  \item \lean{RowSamplingLeverageComputedBasis.lean}: actual-input
  \(A=UC\) equation (7) endpoint and stored-basis table equation (7) endpoint
  for the three rounded-copy modes
  \(\fl_{\rm mul}(U_{ij},1)\), \(\fl_{\rm add}(U_{ij},0)\), and
  \(\fl_{\rm sub}(U_{ij},0)\).
  \item \lean{MatrixAlgebra.lean}: Frobenius norm, vector 2-norm, and the
  Frobenius-to-operator bound used in equation (7).
  \item \lean{DotProduct.lean}: floating-point dot-product stability reused by
  the fully floating-point Gram theorem.
  \item \lean{FiniteProbability.lean}: finite Markov, union-bound,
  monotonicity, and expectation lemmas.
\end{itemize}

\end{document}
